\chapter{容器}
\label{ch:11}

容器的选择取决于怎样访问数据：按固定下标访问、持续追加、用编号查找，还是只判断是否已出现。数组、List、Map、Set 分别给这些任务不同的结构与约定，不能只按写法长短选择。

本章用二维数组观察边界与复制，再通过 Any 容器区分“允许多种内容”与“读取时不需要类型检查”。之后把长度、容量与迭代期间的修改联系起来，为第 17 章的遍历状态做准备。

\section{固定数组与嵌套布局}

有些数据天然有固定形状，例如两行三列的表格。数组把这种形状与元素类型一起记录下来，访问时每个下标分别选择一层位置。先明确行和列，再计算下标，比把它想成任意长度列表更容易检查边界。

下面用数字编码每格位置，再复制整张表。语言数组的布局不能直接替代 C ABI 连续缓冲区约定；checked 的运行期检查也与静态类型、借用检查分别承担职责。

下面初始化二维数组，再复制并修改副本。每个下标分别受对应维度约束。

\par
\medskip
\begin{lstlisting}[style=chali]
void main()
{
    int64 matrix[2][3];
    for (int64 row in range(0, 2))
    {
        for (int64 column in range(0, 3))
        {
            matrix[row][column] = row * 10 + column;
        }
    }
    println("{} {}", matrix[0][2], matrix[1][2]);
    int64 copied[2][3] = matrix;
    copied[1][2] = 99;
    println("{} {}", matrix[1][2], copied[1][2]);
}
\end{lstlisting}
\par
\medskip

\par
\medskip
\begin{lstlisting}[style=out]
2 12
12 99
\end{lstlisting}
\par
\medskip

\subsection*{代码详解}

matrix 的两个维度分别为 2 行和 3 列，下标均从 0 开始。外层 range 选择行，内层选择列，公式 \texttt{row * 10 + column} 让每格的来源易于辨认；第一行是 0、1、2，第二行是 10、11、12，因此首次输出为 2 和 12。

copied 使用同形状数组接收 matrix 的值，再把副本右下角改为 99。第二次输出原值 12 与副本 99，展示二维数组也遵守独立修改规则。维度控制合法下标，不能把 \texttt{[\allowbreak{}2]\allowbreak{}[\allowbreak{}3]\allowbreak{}} 当成可访问的最后一个元素；末元素下标实际是 \texttt{[\allowbreak{}1]\allowbreak{}[\allowbreak{}2]\allowbreak{}}。

\section{按需求选择 List、Map 和 Set}

List 保留序列位置，Map 保存键与值，Set 保存不重复的键。Map 与 Set 的键需要满足哈希和相等协议。Any 可以容纳不同静态类型，但读取时需要恢复或匹配具体类型；已知类型时优先使用具体泛型实参。

\par
\medskip
\begin{lstlisting}[style=chali]
void main()
{
    List<Any> values = [1, "two", true];
    values.push(1.5);
    Any text = values[1];
    println("{} {}", values.len(), tryGet<String>(&text).unwrap());

    Set<Any> unique;
    Any key = "key";
    unique.add(key);
    unique.add(key);
    unique.add(2);
    println("{}", unique.len());

    Map<Any, Any> mixed;
    mixed.put(key, "stored");
    mixed.put(3, true);
    Any saved = mixed[key];
    println("{} {}", mixed.len(), tryGet<String>(&saved).unwrap());

    Map<String, Any> record = {"count": 3, "name": "ChaLi"};
    Any count = record["count"];
    println("{}", tryGet<int64>(&count).unwrap());
    saved.clear();
    println("{}", saved.isEmpty());
}
\end{lstlisting}
\par
\medskip

\par
\medskip
\begin{lstlisting}[style=out]
4 two
2
2 stored
3
1
\end{lstlisting}
\par
\medskip

\subsection*{代码详解}

List<Any> 允许每个元素保存不同具体类型，但取出时仍得到 Any，需要 \texttt{try\allowbreak{}Get<\allowbreak{}String>\allowbreak{}} 等类型化查询确认内容。最初三项加上 1.5 后长度为 4；下标 1 的元素确为字符串，所以本例可 unwrap 得到 two。若查询类型不同，就应处理失败，不能盲目 unwrap。

Set 对重复 key 去重，因此两次加入同一字符串和一次整数后长度为 2。Map 按键保存值，混合键的映射有两项，查询 key 得到 stored。record 的键统一为 String，值仍为 Any，count 需按 int64 取出。最后 clear 只清空局部 saved 中的动态值，isEmpty 变真；它不是清空整个 mixed 映射。

\section{容量、长度与清理}

长度是有效元素数，容量是已分配空间能容纳的数量。追加时先检查 len + need 的溢出和上界；空间不足再按几何增长满足本次需求。删除、clear 和缩短长度保留缓冲容量，但被移除的拥有型元素仍需清理。扩容可以回收旧块，drop 释放最终缓冲区。

最大容量还受元素大小、头部、索引、对齐和分配器限制，不能把可表示的整数上界当成任何元素类型都能实际分配的数量。

\section{借用与遍历}

保存元素引用期间不要做可能移动底层缓冲区的结构修改。复制、移动和迭代器也要区分：复制取得值，移动转移所有权，迭代按协议逐项取得元素。第 17 章用具体例程说明迭代器消费与转换。

\subsection*{小结}

长度表示有效元素数，容量表示可容纳元素的空间，索引表示某个访问位置，三者不能混用。Map 与 Set 还依赖键的相等和哈希关系，Any 则要求取出时确认具体类型。

容器复制遵守值语义，迭代器保存遍历状态，拥有元素又各有清理责任。读容器程序时把这些层次分别标出，便能解释“数据还在但迭代器已结束”或“容量未变但元素已经移出”等现象。

\section{习题}

\begin{enumerate}

\item 分别为成绩序列、学号索引、去重科目选择容器。

\item clear 后容量和元素资源分别怎样？

\item 为什么键的 Hash 和 Equal 必须一致？

\item 列出扩容前需要检查的条件。

\end{enumerate}

\subsection*{参考解答与验收要点}

\begin{enumerate}

\item 通常为 List、Map、Set，仍按实际访问需求决定。

\item 保留容量，清理被移除的拥有型元素。

\item 相等键必须落入一致的索引规则，否则查找可能不一致。

\item 检查加法与倍增溢出、容量上界和分配失败。

\end{enumerate}

下一章：\hyperref[ch:12]{第 12 章}。

\section{番外：刻舟求索引}

改写自《吕氏春秋·察今》刻舟求剑一则。

楚人有涉江者，其剑自列表中坠于第二项，遽契其舟曰：“是吾剑之所从在也。”船工未及言，前一项货物已被删除，后面的项各向前移了一位。

舟止，从其所契者入水求之。得一草鞋。楚人曰：“异哉，剑何以长成了鞋？”船工曰：“剑未变，位置变矣。”楚人怒，谓船工破坏了向后兼容。

船工问：“君所求者，是这第三个位置，抑或那一柄剑？”楚人曰：“当然是剑。若只是第三项，我何必如此伤心？”船工乃给物件各记稳定编号，索引则只管当下排列。

楚人犹不悟，又将一把索引藏进袖中，打算下次扩容后使用。船工叫他停下：“你保存的是位置，不是保鲜袋。”

舟已行矣，而剑之位置不复如故；求剑若此，不亦惑乎？楚人想了半日，终于将船上的“三”改成“重查”。后来有人问他为何刻这两个字，他便指着那只草鞋说：“此乃一次生产事故留下的技术文档。”

\subsection*{技术对照}

索引表示当前排列中的位置，不等于稳定业务标识。插入、删除和结构修改后，旧位置与元素引用的有效性要按具体容器协议判断。
